% !TEX root=main.tex
%%%%%%%%%%%%%%%%%%%%%%%%%%%%%%%%%%%%%%%%%%%%%%%%%%%%%%%%%%%%%%%%%%%%%%
% SECUENCIAS
%
% Strings y Listas
%%%%%%%%%%%%%%%%%%%%%%%%%%%%%%%%%%%%%%%%%%%%%%%%%%%%%%%%%%%%%%%%%%%%%%

\begin{contentbox}{label=Strings}
    \begin{tabular}{C|p{0.5\linewidth}}
        \mintinline{python}|s = "Monty"| & \multirow{2}{=}{Define un objeto de tipo string.} \\
        \mintinline{python}|s = 'Python'| \\
        \mintinline{python}|s[i]| & Accede a la posición \texttt{i} \\
        % \mintinline{python}|s.find(sub)| & Retorna el índice donde empieza \texttt{sub}. \\
        % \mintinline{python}|s.rfind(sub)| & Como \texttt{find}, pero desde la derecha. \\
        % \mintinline{python}|s.index(sub)| & Retorna el índice donde empieza \texttt{sub}, o error si no encuentra.\\
        \mintinline{python}|s.lower()| & Retorna el string en minúscula. \\
        \mintinline{python}|s.upper()| & Retorna el string en mayúscula. \\
        \mintinline{python}|s.strip(t)| & Retorna el string eliminando los caracteres en \texttt{t} de los extremos del string. \\
        \mintinline{python}|s.strip()| & Como \texttt{strip}, pero eliminando espacios en blanco. \\
        % \mintinline{python}|s.lstrip()| & Igual a \texttt{strip}, pero solo a la izquierda. \\
        % \mintinline{python}|s.rstrip()| & Igual a \texttt{strip}, pero solo a la derecha. \\
        % 
        % \mintinline{python}|s.capitalize()| & Retorna el string convertido con el primer carácter a mayúscula, si es una letra, y el resto a minúscula. \\
        % \mintinline{python}|s.title()| & Retorna el string en ``Formato De Título''. \\
        % \mintinline{python}|s.count(sub)| & Retorna la cantidad de apariciones no superpuestas de \texttt{sub}. \\
        \mintinline{python3}|s1 + s2| & Concatena los strings.
    \end{tabular}
\end{contentbox}

\begin{contentbox}{label=Consultas sobre Strings}
    Todos estos métodos devuelven \texttt{True} o \texttt{False}.
    
    \begin{tabular}{C|p{0.5\linewidth}}
        % \mintinline{python}|s.isupper()| & Todas las letras son mayúsculas y hay al menos una. \\
        % \mintinline{python}|s.islower()| & Todas las letras son minúsculas y hay al menos una. \\
        \mintinline{python}|s.isalpha()| & Todos los caracteres son letras. \\
        \mintinline{python}|s.isdigit()| & Todos los caracteres son dígitos. \\
        \mintinline{python}|s.isalnum()| & Todos los caracteres son letras o dígitos. \\
        % \mintinline{python}|s.endswith(t)| & El string termina con \texttt{t}. \\
        % \mintinline{python}|s.startswith(t)| & El string empieza con \texttt{t}.
    \end{tabular}
\end{contentbox}

\begin{contentbox}{label=Índices}
    Los índices de un \textit{string} o lista de largo $n$ parten en 0 y llegan a $n-1$, o parten en -1 (final) y llegan a $-n$ (inicio).
\end{contentbox}

\begin{contentbox}{label=Listas}
    \begin{tabular}{C|p{0.5\linewidth}}
        \mintinline{python}|lista = [1, 2]| & Define un objeto de tipo lista. \\
        \mintinline{python}|lista[i]| & Retorna el elemento en la posición \mintinline{python}|i|. \\
        \mintinline{python}|lista.append(4)| & Añade el elemento \mintinline{python}|4| al final de la lista. \\
        \mintinline{python}|lista[j] = z| & Redefine el valor del elemento en la posición \mintinline{python}|j| de la lista a \mintinline{python}|z|. \\
        % \mintinline{python}|lista.pop(k)| & Retorna el elemento en la posición \mintinline{python}|k| y lo elimina de \mintinline{python}|lista|. Sin parámetros retorna y elimina el último elemento. \\
        \mintinline{python}|lista.pop(k)| & Saca de la lista y retorna el \mintinline{python}|k|-ésimo elemento \\
        % \mintinline{python}|lista.count(c)| & Retorna la cantidad de apariciones del elemento \mintinline{python}|c|. \\
        % \mintinline{python}|lista.index(d)| & Retorna la posición del elemento \mintinline{python}|d|. \\
        % \mintinline{python}|lista.remove(e)| & Elimina la primera aparición del elemento \mintinline{python}|e|. \\
        % \mintinline{python}|lista.insert(i, f)| & Inserta el elemento \mintinline{python}|f| en la posición \mintinline{python}|i|\footnote{Los elementos siguientes son desplazados en 1 posición a la derecha.}. \\
        \mintinline{python}|lista.sort()| & Ordena \mintinline{python}|lista| en orden creciente. \\
        \mintinline{python}|lista1 + lista2| & Retorna una lista que concatena \mintinline{python}|lista1| y \mintinline{python}|lista2|. \\
        \mintinline{python}|lista * n| & Retorna una lista que concatena \mintinline{python}|lista| \mintinline{python}|n| veces\footnote{\mintinline{python}|n| debe ser entero.}. \\
    \end{tabular}
\end{contentbox}

\begin{contentbox}{label=Largo de una Secuencia}
    El largo de la secuencia \mintinline{python}|seq| se obtiene mediante la función nativa \mintinline{python}|len(seq)|. Se entiende por ``largo'' de una secuencia como el total de elementos que hay en esta.
    
    En el caso de una lista, si tiene como elemento otra lista, este elemento sigue contando como uno solo.
\end{contentbox}

\begin{contentbox}{label=String a Lista y Viceversa}
    \begin{itemize}
        \item \mintinline{python}|s.split(sep)| separa el string utilizando \texttt{sep} como separador y retorna una lista cuyos elementos son los fragmentos del string.
        \item \mintinline{python}|sep.join(lista)| une los elementos de \texttt{lista}, separados por \texttt{sep}, en un único string. Solo funciona si los elementos de la lista son strings.
    \end{itemize}
\end{contentbox}

\begin{contentbox}{label=\textit{Slices} y Copias}
    % Pueden accederse \alert{cortes} (\textit{slices}) de una secuencia como una lista o string mediante la \alert{notación slice}. Su notación básica es
    Un \alert{corte} (\textit{slice}) de una secuencia es una \textit{subsecuencia} de la original:
    \begin{center}
        \mintinline{python}|objeto[a:b:c]|
    \end{center}
    Inicia en \mintinline{python}|a|, llega a \mintinline{python}|b| (sin incluirlo) y salta de \mintinline{python}|c| en \mintinline{python}|c| (entero positivo o negativo).
    
    Esta notación tiene las siguientes variaciones:
    \begin{itemize}
        \item Si se omite \mintinline{python}|a| (pero no \mintinline{python}|:|), inicia desde el comienzo.
        \item Si se omite \mintinline{python}|b| (pero no \mintinline{python}|:|), llega hasta el final.
        \item Si se omite \mintinline{python}|c| (incluyendo o no \mintinline{python}|:|), se asume 1.
    \end{itemize}
\end{contentbox}
